\chapter{Results} \label{results}

This chapter presents the main findings from the study described in Chapter \ref{methodology}. The results are presented as follows: first, the participants' profile is described. Translation and PE times are then presented and compared, followed by the analysis of the screen activity tracking data. Process data are complemented with the interviews concerning the participants' experience and impressions of the tasks, and their most relevant comments are summarised. Finally, the analysis of the specific GFL strategies used and the success of their application is presented and gender bias in the MT outputs is briefly discussed.

%----------------------------------------------------------------------------------------
%	PARTICIPANTS
%----------------------------------------------------------------------------------------

\section{Participants' profile}
As initially intended, 12 participants were recruited for the translation experiment. Six received the translation assignment (Task: TR in Table \ref{tab:dem}) and six the PE one (Task: PE in Table \ref{tab:dem}). Participants were aged from 25 to 59 years old. Eight identified as women, three as men, and one as non-binary. Their work experience spanned from 3--5 to more than 20 years, and all except for one had some experience with PE. For the PE assignment, participants with extensive (n=4) or little (n=2) experience in this task were selected. In addition, all participants indicated that they used GFL in their daily work except for a patent translator who stated that this is not desired in the field. The majority of the participants (n=7) indicated that they normally used gender-inclusive characters. Three indicated an alternation with gender-neutral rewording based on the client's requirements and/or the assignment whereas one used all strategies. An overview of the participants' profiles is given in Table \ref{tab:dem}. Their reasons for using GFL include (i) the desire to include everyone in language, (ii) to keep up with linguistic change, and (iii) feminist beliefs. 

\begin{table*}
    \caption{Participants' profiles}
    \resizebox{\textwidth}{!}{\begin{tabular}{llllllll}
    \lsptoprule
        \textbf{Participant} & \textbf{Task} & \textbf{Age} & \textbf{Gender} & \textbf{Work Experience} & \textbf{PE Experience} & \textbf{GFL Experience} & \textbf{GFL Use} \\
        \midrule
        P1 & TR & 39-45 & Woman & 11-15 & Little & Little & Depends on client/assignment \\
        P2 & TR & 25-31 & Woman & 3-5 & Very Little & Yes & Gender-inclusive characters \\
        P3 & TR & 25-31 & Man & 3-5 & No & Little & Gender-inclusive characters \\
        P4 & TR & 32-38 & Non-Binary & 6-10 & Little & Yes & All \\
        P5 & TR & 32-38 & Woman & 6-10 & Very Little & Yes & Gender-inclusive characters \\
        P6 & TR & 32-38 & Woman & 6-10 & Very Little & Little & Gender-inclusive characters \\

        \tablevspace
        P7 & PE & 32-38 & Man & 6-10 & Extensive & Yes & Depends on client/assignment \\
        P8 & PE & 39-45 & Woman & 16-20 & Extensive & Little & No \\
        P9 & PE & 53-59 & Woman & 20+ & Little & Little & Gender-inclusive characters \\
        P10 & PE & 32-38 & Woman & 6-10 & Extensive & Yes & Depends on client/assignment \\
        P11 & PE & 25-31 & Woman & 3-5 & Little & Little & Gender-inclusive characters \\
        P12 & PE & 39-45 & Man & 11-15 & Extensive & Yes & Gender-inclusive characters \\
        \lspbottomrule
    \end{tabular}}
    \label{tab:dem}
\end{table*}

Participants, with the exception of P8 who did not use GFL in their daily work, were also asked to rate the difficulty of GFL on a Likert scale from 1 to 5 where one stands for ``very difficult'' and five for ``very easy''. The majority (n=4) indicated ``neutral'' and the rest were divided between ``easy'' (n=3) and ``difficult'' (n=4), as shown in Figure \ref{fig:gfldifficulty}, where the percentages relate to the number of participants giving this answer. Only one participant indicated ``very easy''. Finally, participants were also asked to comment on challenging and simple aspects of GFL use. Challenges included: (i) the use of neopronouns, especially because they are not yet established in the German language; (ii) the need to adapt and select strategies based on context; (iii) the selection of gender-neutral alternatives to gendered terms; (iv) consistent use of the strategies; and (v) negative impact on the readability of target texts. The use of gender-inclusive characters is simple according to two participants whereas two indicated that GFL is simple regardless of the strategy. 

\begin{figure}[ht]
    \centering
    \includegraphics[width=0.8\textwidth]{figures/GFLDifficulty.png}
    \caption{Ratings of GFL difficulty}
    \label{fig:gfldifficulty}
\end{figure}

%----------------------------------------------------------------------------------------
%	TEMPORAL EFFORT
%----------------------------------------------------------------------------------------

\section{Translation and PE times}
Figure \ref{fig:generaltimes} shows the translation and PE times for each participant in each assignment. It must be noted here that the assignment order was not changed among participants and the most challenging approach to GFL, i.e. neosystems, was applied in the third text. Longer translation and PE times could hence be caused not only by the specific strategy applied but also by participants' tiredness.

First, a difference was found between tasks. PE was generally faster than translation in the first two assignments, i.e. gender-neutral rewording and gender-inclusive characters. Second, differences in speed among participants were considerable. For instance, Figure \ref{fig:generaltimes} clearly shows that P2, and, in part, P10, spent longer times on the third assignment as opposed to the other participants. As observed in \ref{sec:technicaleffort}, they did not split their screen so as to have the translation open alongside the GFL handout, a common strategy among their study colleagues. Finally, neosystems required the longest and gender-inclusive characters the shortest processing time. These differences, however, were not statistically significant. The main findings from the analysis of the translation and PE times are presented in more detail below.

\begin{figure}[ht]
    \centering
    \includegraphics[width=0.9\textwidth]{figures/GeneralTimes.pdf}
    \caption{Translation and PE times for each assignment in minutes}
    \label{fig:generaltimes}
\end{figure}


Table \ref{tab:times} reports the median translation and PE times for each assignment along with the standard deviation (SD). SD values are reported to show the variability in translation and PE speed among participants; this prompted the use of the median in this chapter as a better indicator of the data's central tendency. The translation and PE of the second text with gender-inclusive characters required the least time, i.e. a median of 00:21:20 minutes. The median value was higher in the first assignment, i.e. 00:24:56 minutes. Finally, neosystems took the longest median time, i.e. 00:27:02 minutes.  

\begin{table}[ht]
    \centering
    \caption{Overview of translation and PE times among assignments}
    \begin{tabular}{ccc} \lsptoprule
        \textbf{Assignment} & \textbf{Median} & \textbf{SD} \\
\midrule
        \textbf{1} & 00:24:56 & 00:06:36 \\
        \textbf{2} & 00:21:20 & 00:07:33 \\
        \textbf{3} & 00:27:02 & 00:09:06 \\
\lspbottomrule
    \end{tabular}
    \label{tab:times}
\end{table}

To test whether the task affected the participants' speed, a Welch's test was performed on the rank-transformed data. This test was selected since no equal variance and normal distribution of the data could be assumed. A statistically significant difference between translation and PE was found in the first (\textit{t}[9.98]=1.98, $p$-value=0.038, \textit{d}=1.14) and in the second assignment (\textit{t}[9.37]=3.07, $p$-value=0.006, \textit{d}=1.77). The difference in the third assignment was not statistically significant (\textit{t}[10.0]=0.307, $p$-value=0.383, \textit{d}=0.177). The result suggests that PE was faster than translation in the first and second assignment, but not in the third. These differences can also be seen in Figure \ref{fig:boxplots}, where the box plots compare temporal differences between the translation (TR) and PE (PE) task among assignments.

\begin{figure}[ht]
    \centering
    \includegraphics[width=1.0\textwidth]{figures/boxplots.pdf}
    \caption{Box plots showing differences in translation and PE times (seconds) among assignments}
    \label{fig:boxplots}
\end{figure}

As shown in Table \ref{tab:times2}, the translation task took longer than PE in each assignment. Participants in the translation group completed the first assignment in a median of 00:28:11 minutes -- about eight minutes longer than participants in the PE group with 00:19:58 minutes. A similar difference was found when using gender-inclusive characters: translation required a median of 00:27:48 minutes and PE, 00:20:12. In the case of neosystems, the difference among tasks increased: translation took a median of 00:33:21 minutes and PE, 00:21:00. 

Differences in translation times among assignments were considerable. Gender-neutral rewording and gender-inclusive characters took a median of 00:28:11 minutes and 00:27:48 minutes respectively, increasing by five minutes when using neosystems, i.e. 00:33:21. Median values for PE times were similar among assignments: the first required 00:19:58 minutes, the second 00:20:12, and the third 00:21:00 minutes.

\begin{table}[ht]
    \caption{Comparison of translation and PE times among assignments}
    \centering
    \begin{tabular}{llrr} \lsptoprule
    \textbf{Assignment} & \textbf{Descriptive statistics} & \textbf{TR} & \textbf{PE} \\
\midrule
    1 & Mean & 00:27:45 & 00:19:59 \\
      & Median & 00:28:11 & 00:19:58 \\
      & Shortest Time & 00:21:16 & 00:10:33 \\
      & Longest Time & 00:34:00 & 00:26:12 \\
      & SD & 00:04:49 & 00:06:03 \\
      \midrule
    2 & Mean & 00:27:53 & 00:17:49 \\
      & Median & 00:27:48 & 00:20:12 \\
      & Shortest Time & 00:19:59 & 00:09:26 \\
      & Longest Time & 00:34:59 & 00:21:53 \\
      & SD & 00:06:23 & 00:04:54 \\
      \midrule
    3 & Mean & 00:29:54 & 00:24:04 \\
      & Median & 00:33:21 & 00:21:00 \\
      & Shortest Time & 00:18:02 & 00:12:32 \\
      & Longest Time & 00:40:04 & 00:35:32 \\
      & SD & 00:08:39 & 00:09:59 \\
\lspbottomrule
    \end{tabular}
    \label{tab:times2}
\end{table}

In order to produce a more precise comparison of translation and PE times across assignments, measurements were normalised. A standard approach in TPR is to divide the translation times for each assignment by the number of words in the source texts (\cite{laubli-etal-2013-assessing}), as presented in Table \ref{tab:normalisedtimes}. The results are separated for translation and PE because the translation times were divided by the number of words in the English-language source texts, but PE times by the number of words in the MT outputs. Again, the median was used, but the SD is also reported. To test whether the assignments had an impact on translation and PE times, a rANOVA with Greenhouse-Geisser correction was computed. The test did not show a statistically significant difference among assignments (\textit{F}[1.59,15.86] =3.046, $p$-value=0.085). SD was generally high and was higher in the third assignment than in the first two and in PE than in translation. This indicates that there were considerable speed differences among participants. 

\begin{table}[ht]
\caption{Median value of translation and PE times, SD (seconds per word), and relative SD by assignment}
    \centering
    \small
    \begin{tabular}{cccc} \lsptoprule
   \textbf{Assignment} & \textbf{TR \& PE Times} & \textbf{TR}  & \textbf{PE}\\
\midrule
   1 & Time (s/word) & $11\pm 1.9$ & $7.9\pm 2.4$\\
     & Relative SD & 17\% & 30\% \\ \midrule
   2 & Time (s/word) & $11.2\pm 2.6$ & $7.9\pm 1.9$ \\
     & Relative SD & 23\% & 24\% \\ \midrule
   3 & Time (s/word) & $13.2\pm 3.4$ & $7.8\pm 3.5$ \\
     & Relative SD & 26\% & 44\% \\
\lspbottomrule
    \end{tabular}
    \label{tab:normalisedtimes}
\end{table}

Table \ref{tab:generaltimes} provides an overview of times per participant for each assignment, with the fastest and slowest assignment for each participant indicated in yellow and in red respectively. The second text with gender-inclusive characters was the fastest assignment for half of the participants, namely P2, P6, P7, P8, P11, and P12. Three participants only, namely P1, P3, and P4, took the longest for this second assignment. Gender-neutral rewording was the fastest for four participants, P1, P5, P9, and P10, and the slowest for three, P6, P7, and P11. Translation and PE times for the third assignment were generally higher and neosystems were the fastest assignment for two participants only, namely P3 and P4. Conversely, P2, P5, P8, P9, P10, and P12 took the longest for this. One important difference between translation and PE emerges: while participants in the translation group were generally faster in the third assignment than in the second, post-editors were generally slower in the third. Perhaps, post-editors felt less under time pressure because of the MT and therefore focused more intensively on the third assignment because neosystems were judged to be more complex.

\begin{table}[ht]
\caption{Participants' translation and PE times per assignment in minutes}
    \centering
    \begin{tabular}{ccccc} \lsptoprule
      \textbf{Participant} & \textbf{Task} & \textbf{1st Assignment} & \textbf{2nd Assignment} & \textbf{3rd Assignment} \\
      \midrule
       P1 & TR & \cellcolor{yellow!25}00:30:38 & \cellcolor{red!25}00:34:34 & 00:34:06 \\
       P2 & TR & 00:30:47 & \cellcolor{yellow!25}00:30:20 & \cellcolor{red!25}00:40:04 \\
       P3 & TR & 00:34:00 & \cellcolor{red!25}00:34:59 & \cellcolor{yellow!25}00:33:44 \\
       P4 & TR & 00:21:16 & \cellcolor{red!25}00:22:12 & \cellcolor{yellow!25}00:18:02 \\
       P5 & TR & \cellcolor{yellow!25}00:24:06 & 00:25:15 & \cellcolor{red!25}00:32:58 \\
       P6 & TR & \cellcolor{red!25}00:25:45 & \cellcolor{yellow!25}00:19:59 & 00:20:32 \\ \midrule
       P7 & PE & \cellcolor{red!25}00:22:12 & \cellcolor{yellow!25}00:20:47 & 00:20:55 \\
       P8 & PE & 00:26:12 & \cellcolor{yellow!25}00:21:53 & \cellcolor{red!25}00:34:43 \\
       P9 & PE & \cellcolor{yellow!25}00:17:09 & 00:20:43 & \cellcolor{red!25}00:21:05 \\
       P10 & PE & \cellcolor{yellow!25}00:17:44 & 00:19:41 & \cellcolor{red!25}00:35:52 \\
       P11 & PE & \cellcolor{red!25}00:26:07 & \cellcolor{yellow!25}00:14:23 & 00:20:44 \\
       P12 & PE & 00:10:33 & \cellcolor{yellow!25}00:09:26 & \cellcolor{red!25}00:11:26 \\
\lspbottomrule
    \end{tabular}
    \label{tab:generaltimes}
\end{table}

As concerns translation speed in general, P4 was generally the fastest participant in the translation group. They were the only non-binary participant in the study and rated the difficulty of using GFL as easy in the pre-study questionnaire. Also, as described in Section \ref{sec:strategiesused}, they selected a neosystem that they already used in their daily life with friends. P4 was the fastest in completing the first and the third assignments. In the second, P6 was the fastest. P6 admittedly had less experience with GFL than P4 and found it difficult to use. Among translators, P3 was generally the slowest participant in the first and second assignments. When using neosystems, P2 needed the longest time. Both rated GFL as neither difficult nor easy, but P2 had more experience with it. P12 was the fastest participant overall, post-editing each text in about ten minutes even though they rated GFL use as difficult. P8 was the slowest participant in the PE group and, in fact, the only participant who did not integrate GFL in their everyday work.

When asked to comment on the impact of GFL on their speed and productivity, half of the participants stated that they felt GFL generally slowed them down. This perception was more common in the translation (n=4) than in the PE group (n=2). Also, five participants (P5, P7, P8, P10, P12) felt the least productive when using neosystems. This impression was confirmed when analysing translation and PE times except for P7. Finally, P4 and P6 felt they were the slowest in the first assignment, which was only confirmed for P6. No participant commented on gender-inclusive characters as the most time-consuming approach to GFL.

To summarise, differences in times could be partially observed between assignments. In translation, gender-neutral rewording and gender-inclusive characters each required a median of 28 minutes, while neosystems took a median of 33 minutes. Neosystems required the longest for half of the participants. PE was significantly faster than translation in the first two assignments and the median for each assignment was about 20 minutes. Furthermore, there were considerable differences in translation and PE among participants, but longer or shorter times did not always correlate with the perceived difficulty of GFL. Finally, the perception of the negative impact of GFL on speed and productivity was more common among translators than post-editors. This might be because PE is usually faster than translation and, as explained in Section \ref{sec:cognitiveeffort}, post-editors felt that the high-quality MT outputs they received improved their productivity.  

%----------------------------------------------------------------------------------------
%	TECHNICAL EFFORT
%----------------------------------------------------------------------------------------
\section{Screen activity tracking data} \label{sec:technicaleffort}
In order to analyse the screen activity and gain insights into the gender-fair translation and PE process, observation protocols were produced with the observation notes and the screen recordings. The observation notes included comments on general approaches to and/or interesting aspects of the process. The screen recordings analysis was focused on the number of activities, i.e. (i) broken units, (ii) changes, (iii) searches, and (iv) pauses, performed by the participants. 

First, as regards general approaches to the translation and PE, the majority of the participants (n=8) looked online for more information concerning the TV series described in the texts and/or the characters and actors mentioned. Only P2, P5, P6, and P12 did not. P3 and P9 also looked online for images of the non-binary actors. P4 even read the biography of Dua Saleh, a celebrity mentioned in the first text, in order to reformulate one text passage that contained the nouns \textit{actor} and \textit{musician}. P4 replaced the former by citing a series in which the actor played (``die auch in Generation mitgespielt hat'', i.e. who also acted in Generation). 

%\begin{figure}[ht]
 %  \includegraphics[width=0.80\textwidth]{figures/screen1.png}
   % \caption{Image search to visualize the non-binary actor mentioned in one text}
    %\label{fig:screen1}
%\end{figure}

Participants also highlighted gender references for which GFL was to be used. P3 and P5 did so in every assignment, i.e. they often translated gender references in the masculine or feminine form, highlighted them as shown in Figure \ref{fig:screen2}, then used GFL in a subsequent step. P8 and P11 adopted the same strategy in the first assignment only, whereas P9 did so for a sentence in the third assignment which they found particularly difficult. P4 used a gender-inclusive character and a neopronoun in the first assignment, putting them into brackets in order to replace them with a gender-neutral strategy in a subsequent step.

A common strategy in the third assignment was to split the screen so as to have the translation open alongside the GFL handout or a similar resource. This allowed P3, P5, P7, P9, P11, and P12 to avoid constantly opening and closing windows with resources on how to use the selected neosystem as P2 and P10 did. Another strategy was to focus on GFL as a separate task from translation and PE. For instance, P9, P11, and P12 corrected gender references first and then post-edited the MT outputs. 

\begin{figure}[ht]
    \centering
    \includegraphics[width=0.80\textwidth]{figures/Screen2.png}
    \caption{Use of the highlight function to later revise GFL}
    \label{fig:screen2}
\end{figure}

Second, the analysis of specific screen activities, i.e. broken units, changes, searches, and pauses, is presented below. In Figure \ref{fig:technical effort}, the median values and sum for each activity are reported for each assignment. 

\begin{figure}[ht]
    \centering
    \includegraphics[width=0.95\textwidth]{figures/technicaleffort.pdf}
    \caption[Overview of the keyboard activity for each assignment]{Overview of the keyboard activity for each assignment with median values}
    \label{fig:technical effort}
\end{figure}

The median number of the observed activities was similar in the first two assignments, i.e. 24.5 and 24 respectively, but considerably increased to a median value of 50 when using neosystems. Activities occurred with different frequencies. Broken units were found sporadically in the observation protocols. Changes and searches were most frequent when participants used neosystems, with a median of 17 and 16.5 occurrences respectively, as against a respective median of 7.5 and 3.5 when gender-neutral rewording, and 13.5 and 2.5 when applying gender-inclusive characters. Pauses were the most frequent when participants reworded in gender-neutral ways, with a median 9.5 times, as against 4.5 in the second and 7 in the third assignment. To test whether the different strategies had an impact on the screen activity tracking, a rANOVA with Greenhouse-Geisser correction was performed. It showed a statistically significant difference in the observed activities among assignments (\textit{F}[1.45,14.46]=10.54, $p$-value =0.003). Post-hoc analysis was conducted with Bonferroni correction to determine the statistical significance of relationships among the specific strategies. There was a statistically significant difference between the first and the third ($p$-value = 0.017) as well as the second and third assignment ($p$-value = 0.021). This means that activity increased in the third assignment and no statistically significant difference was found between the first and the second assignment. 


Differences can also be found in the comparison between translation and PE, as shown in Table \ref{tab:technicaleffort}. The median of screen activities was higher in the PE task in gender-inclusive characters (TR: 20.5 vs PE: 32) and neosystems (TR: 41 vs PE: 52). This difference is largely influenced by the number of searches in the third assignment, which were conducted more frequently in the PE (25.5) than in the translation task (14). In gender-neutral rewording, the analysed screen activities were higher in translation (27) than in PE (16.5). To test whether the task had an influence on activity frequency, a Welch's test was performed on ranked data. The test was selected because the data violated the assumption of equal variance and were not normally distributed. It showed no statistically significant difference between translation and PE in each assignment (first assignment: \textit{t}[7.89]=-1.53, $p$-value=0.166, \textit{d}=-0.881, second assignment: \textit{t}[9.80]=1.06, $p$-value=0.317, \textit{d}=-0.609, third assignment: \textit{t}[9.68]=-0.788, $p$-value=0.450, \textit{d}=-0.455).

\begin{table}[ht]
    \caption[Overview of screen activity tracking with median values by assignment and task]{Overview of screen activity tracking with median values by assignment and task (normalised values based on the number of gendered phrases per assignment in brackets)}
    \centering
    \begin{tabular}{rlrr} \lsptoprule
        \textbf{Assignment} & \textbf{Screen activities} & \textbf{TR} & \textbf{PE} \\
\midrule
        1 & Total & 27 (3) & 16.5 (1.85) \\
          & Broken units & 2 (0.2) & 0.5 (0.05) \\
          & Changes & 6.5 (0.7) & 10 (1.1) \\
          & Searches & 4.5 (0.5) & 0.5 (0.05) \\
          & Pauses & 10 (1.1) & 4.5 (0.5) \\ \midrule
        2 & Total & 20.5 (1.7) & 32 (2.65) \\
          & Broken units & 0.5 (0.4) & 1 (0.08) \\
          & Changes & 8 (0.65) & 18.5 (1.55) \\
          & Searches & 4 (0.35) & 2 (0.15) \\
          & Pauses & 5 (0.4) & 4.5 (0.38) \\ \midrule
        3 & Total & 41 (4.1) & 52 (5.2) \\
          & Broken units & 1.5 (0.15) & 1.5 (0.15) \\
          & Changes & 12 (1.2) & 18 (1.8) \\
          & Searches & 14 (1.4) & 25.5 (2.55) \\
          & Pauses & 7.5 (0.75) & 2.5 (0.25) \\
\lspbottomrule
    \end{tabular}
    \label{tab:technicaleffort}
\end{table}

Post-editors made more changes, whereas translators made more pauses. A higher number of changes in the PE task was expected since participants had to correct almost every machine-translated gender reference, while participants in the translation task could produce an appropriate solution without modifying it. Nevertheless, there were differences among assignments in the number of changes. While translating, the median for changes was 6.5, 8, and 12 respectively. Hence, translators made more changes when using neosystems. While in PE, the median for changes was 10, 18, and 18.5 respectively, hence there were no difference between gender-inclusive characters and neosystems. Pauses were more frequent in the translation task in the first (TR: 10 vs PE: 4.5) and third (TR: 7.5 vs PE: 2.5) assignments. 

Since the number of changes cannot be directly compared between tasks (i.e. translation and PE), the analysis of the screen recordings was repeated and the count of changes was omitted as shown in Table \ref{tab:technicaleffort_nochanges}. In this case too, the highest screen activity was found in the third assignment both in translation and in PE. PE in this assignment required a higher number of activities than translation. Nevertheless, the difference seems to be caused by two participants in particular, as shown in Table \ref{tab:technicaleffortt_comparisonParticipant}. P8 and P10 did 71 and 72 searches respectively. When using gender-neutral rewording, the frequency of the observed activities was higher in the translation task, whereas no differences were found in the second text with gender-inclusive characters. 

\begin{table}[ht]
    \caption{Total number of screen activities with their median in brackets after omitting the count of changes}
    \centering
    \begin{tabular}{rrrr}
    \lsptoprule
       \textbf{Assignment} & \textbf{TR} & \textbf{PE} & \textbf{Total} \\
       \midrule
        1 & 103(18) & 66(10) & 169(13.5) \\ 
        2 & 77(9.5) & 74(9.5) & 151(10) \\ 
        3 & 147(27) & 246(35) & 393(30.5) \\
\lspbottomrule
    \end{tabular}
    \label{tab:technicaleffort_nochanges}
\end{table}

Table \ref{tab:technicaleffortt_comparisonParticipant} shows the total number of observed activities per participant for each assignment, where B indicates broken units, C changes, S searches, P pauses, and T the total of their sum. For each participant, the assignment with the lowest and highest screen activity is highlighted in yellow and red respectively. The table reflects the intuitive notion that screen activity was the highest in the third assignment (using neosystems) for ten participants, i.e. P1, P2, P3, P5, P7, P8, P9, P10, P11, and P12. For instance, an increased number of changes and searches was observed. The lowest activity was observed in the first assignment for six participants, namely P1, P5, P7, P9, P10 and P12, and in the second for five, i.e. P2, P3, P6, P8, P11.


\begin{table}[ht]
\caption{Comparison of screen activity tracking per participant in each assignment}
    \centering 
    \resizebox{\textwidth}{!}{\begin{tabular}{rrrrrrrrrrrrrrrr}
    \lsptoprule
        P& \multicolumn{5}{c}{\textbf{1st Ass.}}& \multicolumn{5}{c}{\textbf{2nd Ass.}}& \multicolumn{5}{c}{\textbf{3rd Ass.}}\\
\midrule
 & B& C& S& P& T& B& C& S& P& T& B& C& S& P&T\\
       \textbf{P1} & 1 & 6 & 12 & 10 &\cellcolor{yellow!25}29& 0 & 6 & 11 & 13 &\cellcolor{red!25}30& 0 & 6 & 17 & 7 &\cellcolor{red!25}30\\
       \textbf{P2} & 3 & 6 & 5 & 10 &24& 0 & 11 & 7 & 3 &\cellcolor{yellow!25}21& 6 & 16 & 16 & 14 &\cellcolor{red!25}52\\
       \textbf{P3} & 0 & 17 & 11 & 2 &30& 0 & 6 & 2 & 1 &\cellcolor{yellow!25}9& 2 & 30 & 21 & 8 &\cellcolor{red!25}61\\
       \textbf{P4} & 6 & 14 & 1 & 12 &\cellcolor{red!25}33& 1 & 10 & 6 & 3 &20& 1 & 8 & 1 & 5 &\cellcolor{yellow!25}15\\
       \textbf{P5} & 0 & 6 & 3 & 9 &\cellcolor{yellow!25}18& 3 & 17 & 1 & 16 &37& 2 & 22 & 12 & 16 &\cellcolor{red!25}52\\
       \textbf{P6} & 4& 7 & 4 & 10 &\cellcolor{red!25}25& 2 & 6 & 1 & 7 &\cellcolor{yellow!25}16& 1 & 5 & 11 & 7 &24\\
\midrule
       \textbf{P7} & 0 & 8 & 0 & 1 &\cellcolor{yellow!25}9& 0 & 12 & 0 & 0 & 12 & 1 & 5 & 11 & 7 &\cellcolor{red!25}18\\
       \textbf{P8} & 4 & 19 & 7 & 23 & 53 & 1 & 21 & 21 & 8 &\cellcolor{yellow!25}51 & 4 & 27 & 72 & 9 &\cellcolor{red!25}112\\
       \textbf{P9} & 0 & 7 & 0 & 3 &\cellcolor{yellow!25}10 & 2 & 25 & 3 & 7 & 37 & 0 & 16 & 20 & 12 &\cellcolor{red!25}48\\
       \textbf{P10} & 0 & 7 & 4 & 6 &\cellcolor{yellow!25}17 & 1 & 15 & 16 & 5 & 37 & 3 & 20 & 71 & 1 &\cellcolor{red!25}95\\
       \textbf{P11} & 2 & 14 & 1 & 11 &28& 2 & 20 & 1 & 4 &\cellcolor{yellow!25}27& 4 & 18 & 31 & 3 &\cellcolor{red!25}56\\
       \textbf{P12} & 1 & 12 & 0 & 3 &\cellcolor{yellow!25}16& 1 & 17 & 0 & 2 &20& 0 & 19 & 7 & 2 &\cellcolor{red!25}28\\
\lspbottomrule
    \end{tabular}}
    \label{tab:technicaleffortt_comparisonParticipant}
\end{table}

P7 was the participant for whom the lowest number of activities was observed, with a total of 9, 12, and 18 activities respectively for the three assignments. They made no searches and almost no pauses in the first two assignments. They were also the participant who performed the lowest number of activities in the first assignment, followed by P9. The lowest screen activity in the second assignment was registered by P3, followed by P7. When using neosystems, P4 was the participant who produced the least screen activity tracking data, with e.g. only one search. 

P8 was the participant whose screen activity was the highest in each assignment with a considerable number of changes, i.e. 19, 21, and 27 respectively. This indicates that they were neither satisfied with nor sure about the use of GFL. They also made several pauses (23) in the first assignment while thinking of possible gender-neutral solutions and did a very high number of searches (72) in the third. The difference in the number of searches by P8 and P10 (72 and 71 respectively) versus the other participants is substantial, with P11 conducting 31 searches and the others even fewer.

In Table \ref{tab:temporal/technicaleffort}, translation and PE times are compared with the total number of analysed screen activities. The lowest values are highlighted in yellow and the highest in red. For half of the participants (P2, P5, P6, P8, P9, and P10), there was a correspondence between longer translation times and higher activity. For P3, conversely, longer translation times corresponded to lower activity. A higher number of screen activities was observed for P7 and P11 when using neosystems, while for the first assignment with gender-neutral rewording, they needed more time. A Kendall's Tau was performed to test for a relationship between translation/PE times and screen activity; this showed no correlation ($\tau$ = 0.074, $p$-value = 0.530). 


\begin{table}[ht]
\caption{Triangulation of translation and PE times with screen activity tracking data}
    \centering
    \begin{tabular}{crrr} \lsptoprule
      \textbf{Participant} & \textbf{1st Assignment} & \textbf{2nd Assignment} & \textbf{3rd Assignment} \\
\midrule
       P1 & \cellcolor{yellow!25}00:30:38 & \cellcolor{red!25}00:34:34 & 00:34:06 \\
          & \cellcolor{yellow!25}29 & \cellcolor{red!25}30 & \cellcolor{red!25}30 \\
       P2 & 00:30:47 & \cellcolor{yellow!25}00:30:20 & \cellcolor{red!25}00:40:04 \\
          & 24 & \cellcolor{yellow!25}21 & \cellcolor{red!25}52 \\
       P3 & 00:34:00 & \cellcolor{red!25}00:34:59 & \cellcolor{yellow!25}00:33:44 \\
          & 30 & \cellcolor{yellow!25}9 & \cellcolor{red!25}61  \\
       P4 & 00:21:16 & \cellcolor{red!25}00:22:12 & \cellcolor{yellow!25}00:18:02 \\
          & \cellcolor{red!25}33 & 20 & \cellcolor{yellow!25}15 \\
       P5 & \cellcolor{yellow!25}00:24:06 & 00:25:15 & \cellcolor{red!25}00:32:58 \\
          & \cellcolor{yellow!25}18 & 37 & \cellcolor{red!25}52 \\
       P6 & \cellcolor{red!25}00:25:45 & \cellcolor{yellow!25}00:19:59 & 00:20:32 \\ 
          & \cellcolor{red!25}25 & \cellcolor{yellow!25}16 & 24 \\
\midrule
       P7 & \cellcolor{red!25}00:22:12 & \cellcolor{yellow!25}00:20:47 & 00:20:55 \\
          & \cellcolor{yellow!25}9 & 12 & \cellcolor{red!25}18 \\
       P8 & 00:26:12 & \cellcolor{yellow!25}00:21:53 & \cellcolor{red!25}00:34:43 \\
          & 53 & \cellcolor{yellow!25}51 & \cellcolor{red!25}112 \\
       P9 & \cellcolor{yellow!25}00:17:09 & 00:20:43 & \cellcolor{red!25}00:21:05 \\
          & \cellcolor{yellow!25}10 & 37 & \cellcolor{red!25}48 \\
       P10 & \cellcolor{yellow!25}00:17:44 & 00:19:41 & \cellcolor{red!25}00:35:52 \\
           & \cellcolor{yellow!25}17 & 37 & \cellcolor{red!25}95 \\
       P11 & \cellcolor{red!25}00:26:07 & \cellcolor{yellow!25}00:14:23 & 00:20:44 \\
           & 28 & \cellcolor{yellow!25}27 & \cellcolor{red!25}56\\
       P12 & 00:10:33 & \cellcolor{yellow!25}00:09:26 & \cellcolor{red!25}00:11:26 \\
           & \cellcolor{yellow!25}16 & 20 & \cellcolor{red!25}28\\
\lspbottomrule
    \end{tabular}
    \label{tab:temporal/technicaleffort}
\end{table}

Higher screen activity does not correlate with the perceived difficulty of GFL either. P10 indicated that GFL use is very easy but their screen activity was one of the highest during the study. P12 indicated that GFL use is difficult, but they were the participant who produced the least screen activity tracking data. Use, or actually lack thereof, in everyday work-life seems to impact the number of screen activities. P8 was the sole participant who did not use GFL in their work and their screen activity was the highest in the study. 

In order to shed light on which gender references require more effort in terms of screen activity in translation and PE, the three texts translated during the experiment were divided into segments based on punctuation rules. For each, the number of activities performed was annotated. Also, notes were taken to reconstruct the translation and PE process of GFL. An overview of the total activities performed in each text segment is presented below, along with examples of gender-specific terms that were particularly challenging in the source and target texts.  

The first text contained the following seven segments (analysed gender references are shown in bold):
\begin{enumerate}
\setlength\itemsep{-0.5em}
    \item Sex Education on Netflix is known for its no-holds-barred approach to sex and sexuality.
    \item In season three, the show takes a similar approach to \textbf{gender identity}, shining a bright light on \textbf{those who identify as non-binary}.
    \item Leading the way is \textbf{a non-binary actor and musician} Dua Saleh as Cal, \textbf{a new student} at Moordale who clashes directly with Headmistress Hope (Jemima Kirke).
    \item Hope is determined to turn Moordale's reputation as a ``sex school'' around and in doing so, she sets out to squash \textbf{students'} sexual freedoms and \textbf{their} right to express \textbf{their} identity.
    \item Cal stands up against Hope, refusing to wear the new Moordale school uniform that is deemed acceptable in Hope's eyes. 
    \item \textbf{They} refuse to wear the girls [sic] uniform and is [sic] reluctant to wear the tight-fighting [sic] boys [sic] uniform. 
    \item Cal likes to keep \textbf{their} uniform loose and baggy, as it makes \textbf{them} feel more comfortable in \textbf{who they} are, but Hope just won't listen.
\end{enumerate}


%da qui ricominciare con analisi
As shown in Table \ref{tab:technicaleffort_segment_first}, no activity was performed in the first and fifth segments. The third showed the highest tracking activity, e.g. 50 changes and 47 pauses (values highlighted in red). In the last segment as well, participants faced great difficulties. The greatest problems thus relate to the translation of ``actor and musician'', ``student'', and the singular \textit{they} pronoun in different cases. 

\begin{table}[ht]
    \caption[Total of screen events per segment in the first assignment]{Total for screen events per segment in the first assignment (Translation|PE)}
    \centering
    \begin{tabular}{cccccc} \lsptoprule
        \textbf{Segment} & \textbf{Broken units} & \textbf{Changes} & \textbf{Searches} & \textbf{Pauses} & \textbf{Total}\\
\midrule
        \textbf{1} & 0|0 & 0|0 & 0|0 & 0|0 & 0|0 \\
        \textbf{2} & 1|0 & 4|4 & 8|2 & 0|0 & 13|6 \\
        \textbf{3} & 3|2 & \cellcolor{red!25}21|29 & \cellcolor{red!25}18|4 & \cellcolor{red!25}25|22 & \cellcolor{red!25}67|57 \\
        \textbf{4} & 2|1 & 5|8 & 2|1 & 6|3 & 15|13 \\
        \textbf{5} & 0|0 & 2|0 & 0|0 & 0|0 & 2|0 \\
        \textbf{6} & 1|2 & 6|12 & 5|2 & 6|10 & 18|24 \\
        \textbf{7} & 7|2 & \cellcolor{red!25}18|14 & 3|3 & \cellcolor{red!25}16|12 & \cellcolor{red!25}44|31 \\
\lspbottomrule
    \end{tabular}
    \label{tab:technicaleffort_segment_first}
\end{table}

For instance, before translating ``student'' P1 made a ten-second pause. They used the word ``Neuzugang'' (new entry) and made another pause, this time of nine seconds. They then looked on the internet for the frequency of use of ``Neuzugang'' and ``Neuankömmling'' (newcomer) and made another ten-second pause. Then, they deleted the former term and replaced it with the latter. They made a 15-second pause, deleted the term, made a 6-second pause and finally opted for ``neuen Mitglied'' (new member). 

While translating the third segment, P2 looked four times in the provided handout for inspiration. Also, before finding a solution for ``actor and musician'', they made two long pauses of 38 and ten seconds respectively. Then, they started writing ``Schau'' (act), made another pause of nine seconds, wrote ``Schauspie'' (acting, incomplete word), deleted it, and wrote ``Schauspielperson und Musik'' (acting person and music). They did not complete the translation but made a 21-second pause. Finally, they opted for ``Musikmensch'' (music person). 

P8 had difficulties in translating singular \textit{they} and made a long pause of 90 seconds in the seventh segment, before PE the MT for ``it makes them feel more comfortable in who they are''. They also looked three times online for synonyms for ``comfortable'' in order to restructure the sentence and made six changes to the segment. Finally, during the revision phase, they made another seven pauses lasting from 15 to 40 seconds.

The second text contained six segments (analysed gender references are shown in bold):
\begin{enumerate}
\setlength\itemsep{-0.5em}
    \item Grey’s Anatomy has \textbf{its first non-binary doctor} after E. R. Fightmaster was promoted from a \textbf{recurring guest star} to a \textbf{recurring cast member}.
    \item According to Variety, \textbf{the actor} has landed a recurring role as \textbf{Dr.} Kai Bartley after making \textbf{their} first appearance in Season 18.
    \item \textbf{Viewers} were introduced to \textbf{them} in episode three when Meredith Grey (Ellen Pompeo) was contacted by a well-funded Minnesota hospital that wants her to join their mission to cure Parkinson’s Disease. 
    \item While \textbf{audiences} have not seen much of \textbf{the brilliant neuroscientist}, played by \textbf{nonbinary actor} E.R. Fightmaster, \textbf{they} are already eager for more. 
    \item ‘\textbf{They}  are dedicated to \textbf{their}  craft and extremely talented at what \textbf{they} do. 
    \item Confident as hell and able to make even the most detailed and mundane science seem exciting and cool, Kai and Amelia bond over \textbf{their} shared love of medicine and the brain,’ the ABC character description for Kai reads.
\end{enumerate}

During the translation process, one participant looked at the handout before starting their revision. Other participants generally used resources when they encountered a term for which the translation or the gender-fair alternative was difficult. As shown in Table \ref{tab:technicaleffort_segment_second}, the first and the fourth segments caused difficulties for both translators and post-editors (values in red). The screen activity was high in the second and third segments as well. Terms that caused difficulties were ``doctor'', ``recurring guest star'', ``recurring cast member'', and singular \textit{they}. In addition, the translation of the fourth segment required an abundance of gender-inclusive characters, thus prompting a high number of changes and searches in order to reword the text in an attempt to avoid overusing characters that would impact readability. 

\begin{table}[ht]
    \caption[Total of screen activities per segment in the second assignment]{Total of screen activities per segment in the second assignment (Translation|PE)}
    \centering
    \begin{tabular}{cccccc} \lsptoprule
        \textbf{Segment} & \textbf{Broken units} & \textbf{Changes} & \textbf{Searches} & \textbf{Pauses} & \textbf{Total}\\
\midrule
        \textbf{1} & 5|6 & \cellcolor{red!25}29|26 & \cellcolor{red!25}18|12 & \cellcolor{red!25}20|11 & \cellcolor{red!25}72|55 \\
        \textbf{2} & 1|0 & 9|21 & 2|6 & 7|2 & 19|29 \\
        \textbf{3} & 0|1 & 5|14 & 1|6 & 4|4 & 10|25 \\
        \textbf{4} & 0|0 & \cellcolor{red!25}9|25 & \cellcolor{red!25}5|13 & \cellcolor{red!25}10|3 & \cellcolor{red!25}24|41 \\
        \textbf{5} & 0|0 & 3|17 & 1|0 & 3|4 & 7|21 \\
        \textbf{6} & 0|0 & 1|7 & 0|4 & 1|2 & 2|13 \\
        \textbf{pre-revision} & 0|0 & 0|0 & 1|0 & 0|0 & 1|0 \\
\lspbottomrule
    \end{tabular}
    \label{tab:technicaleffort_segment_second}
\end{table}

P5, for example, was unsure how to translate the term ``doctor'' with gender star because the masculine and feminine forms of the noun differ for the umlaut on the first letter (``Arzt'', ``Ärztin''; male doctor, female doctor). Hence, they first translated the English term using the masculine form, and then highlighted it. Subsequently, they tried to write an alternative, which they immediately deleted. Then they rewrote the sentence, this time using the feminine form. In brackets, they added an alternative translation using the synonym ``Mediziner*in''. During the revision phase, they also searched online how to use gender-inclusive characters in the word ``Arzt''. Similarly, P8 searched online five times how to gender the term. 

P10 faced difficulties while post-editing the MT for ``recurring guest star''. While post-editing, they stopped writing ``wiederkehrende...'' (recurring) and looked online five times for a translation, using different electronic vocabularies. They also made three pauses while PE the segment and opted for rewording the term as ``Star mit wiederkehrenden Gastauftritten'' (star with recurring guest appareances). In this case, however, the use of GFL was not the sole cause of difficulties. Participants struggled to find a German equivalent for the source term. 

Finally, participants made changes because they were unsure regarding the order of masculine and feminine forms when using pronouns and articles. For instance, P11 changed the MT of ``actor'' from ``der Schauspieler'' (masculine article and noun) to ``der*die Schauspieler*in'' (masculine*feminine article +  masculine*feminine noun) and then changed it to ``die*der Schauspieler*in'' (feminine*masculine article + feminine*masculine noun). Similarly, P1 translated using masculine forms first, then read the provided handout and decided to change the order of gendered forms accordingly. 

The third text contained seven segments (analysed gender references are shown in bold):
\begin{enumerate}
\setlength\itemsep{-0.5em}
    \item \textbf{Audiences} at last week's Toronto International Film Festival didn't just watch movies. 
    \item \textbf{They} were also lucky enough to get a sneak peak at what will surely become one of the most celebrated TV series of this fall: Sort Of. 
    \item This is a comedy about \textbf{a gender-fluid nanny} navigating \textbf{their} complicated existence by \textbf{co-creators} Bilal Baig and Fab Filippo. 
    \item The series also features Baig. Baig is \textbf{a South Asian, queer Muslim actor who} was previously known best for \textbf{their} work in theatre, and \textbf{whose} deadpan delivery in Sort Of is one of its many great assets. 
    \item \textbf{They} star as \textbf{nanny} Sabi Mehboob, \textbf{who} is given a chance to move to Berlin in pursuit of something more than \textbf{their} clearly unsatisfying life. 
    \item But after the mother of the children gets in a serious bike accident, Sabi decides to stay put in Toronto. 
    \item As Bilal mentions, Sort Of represented \textbf{their} first major experience with television.
\end{enumerate}
In this case, some participants also read the handout or used other resources before and after their translation and/or PE in order to learn more about how the selected neosystems are used. Table \ref{tab:technicaleffort_segment_third} shows an overview of the screen activity tracking data produced in each segment. The greatest difficulties were faced in the third, fourth, and fifth segments for which the number of changes and searches is particularly high. Difficult gender constituents were ``nanny'', ``co-creators'', adjective declensions, and pronouns. 

\begin{table}[ht]
    \caption[Total of screen activities per segment in the third assignment]{Total screen activities per segment in the third assignment (Translation|PE)}
    \centering
    \begin{tabular}{cccccc} \lsptoprule
        \textbf{Segment} & \textbf{Broken units} & \textbf{Changes} & \textbf{Searches} & \textbf{Pauses} & \textbf{Total}\\
\midrule
        \textbf{pre-translation} & 0|0 & 0|0 & 3|4 & 0|0 & 3|4 \\
        \textbf{1} & 2|0 & 4|1 & 3|1 & 0|0 & 9|2 \\
        \textbf{2} & 0|0 & 1|1 & 3|1 & 2|1 & 6|3 \\
        \textbf{3} & 3|4 & \cellcolor{red!25}33|36 & \cellcolor{red!25}39|84 & \cellcolor{red!25}20|13 & \cellcolor{red!25}95|137 \\
        \textbf{4} & \cellcolor{red!25}5|6 & \cellcolor{red!25}29|35 & \cellcolor{red!25}14|62 & \cellcolor{red!25}18|7 & \cellcolor{red!25}66|110 \\
        \textbf{5} & 2|0 & \cellcolor{red!25}14|28 & \cellcolor{red!25}11|39 & \cellcolor{red!25}15|4 & \cellcolor{red!25}42|71 \\
        \textbf{6} & 0|0 & 3|3 & 2|3 & 1|2 & 6|8 \\
        \textbf{7} & 0|1 & 3|7 & 1|12 & 1|2 & 5|22 \\
        \textbf{pre-revision} & 0|0 & 0|0 & 2|0 & 0|0 & 2|0 \\
\lspbottomrule
    \end{tabular}
    \label{tab:technicaleffort_segment_third}
\end{table}

For instance, P6 looked for a possible translation of ``nanny'' in an online dictionary. Then they searched for ``Kinderfrau'' (synonym for ``nanny'') in the Duden,\footnote{\url{https://www.duden.de}} a German language dictionary. They finally translated the term as ``Kinderpflegens'' (childcare worker with the Ens-form) and then changed it to ``Kinderbetreuens'' (caregiver). P3 struggled with the use of possessives in the Sylvain system. While post-editing the MT for ``whose deadpan delivery'', the participant first opted for omitting the pronoun, then they changed it to ``seiner/ihrer'' (his/her), ``ihrin'' (a tentative use of possessives with the Sylvain system), ``ihrins'', ``seiner/ihrer'', and finally ``ihrin''. They also looked at the handout six times as well as at a paper explaining the system. 

From the observation, it also emerges that most of the participants used the handout or other resources concerning neosystems for the whole duration of the study. During the whole PE of the third assignment, e.g. P8 and P10 conducted 72 and 71 searches respectively. A common strategy was to open the document with source and target text along with the handout or another resource describing the use of the selected neosystem. Almost every time a gender reference was encountered, the resource was consulted. This also resulted in a continually interrupted writing flow, for instance in segment four, where adjectives, nouns, and pronouns occurred with a certain frequency.

To summarise, differences in screen activity tracking data were found among strategies as well as participants. Activity frequency was similar in the first and second assignment, but significantly increased when neosystems were used. When using gender-neutral rewording, participants made more pauses, probably to reflect on gender-neutral alternatives to gender-specific terms. In the third assignment, the number of searches and -- to a lesser extent -- changes increased. There were also differences between translation and PE, but these are less clear-cut. In this case, the number of changes cannot be directly compared. Since machine-translated gender references were often wrong or contained masculine generics, changes in the PE tasks occurred more frequently. However, if this indicator is excluded from the analysis, post-editors performed more activities in the third assignment. Results are probably skewed since P8 and P10 did considerably more searches than the other participants. No difference between translation and PE was found when using gender-inclusive characters, and gender-neutral rewording was the only case where translation required higher activity than PE. This outcome may be coincidental, as inferential statistical analysis indicates that the observed difference lacks statistical significance. Finally, translators made more searches than post-editors in the first two assignments and more pauses in the first and third. While the perceived difficulty of GFL use does not correlate with a higher number of activities, the frequency of its use seems to play an important role. The only participant who does not use GFL in their work, i.e. P8, is the person who made more pauses, changes, and searches.

%----------------------------------------------------------------------------------------
%	COGNITIVE EFFORT
%----------------------------------------------------------------------------------------
\section{Participants' experience} \label{sec:cognitiveeffort}
In order to investigate the participants' personal experience with GFL during the study, retrospective interviews with the translators were conducted at the end of each study session. Participants were asked to comment on difficulties in the source text; on the translation process, with a particular focus on easy and challenging aspects of GFL use; and their impressions of the advantages and disadvantages of using MT  for GFL (see the interview guide and the transcripts in the Appendix). %An overview of the participants' impressions is provided in this section. In contrast to the previous sections, results are presented by immediately disaggregating data by task, i.e. translation and PE, for the sake of brevity.

%insert comparison between the assigments in terms of most difficult
While gender-neutral rewording was the easiest assignment for five participants, a recurring comment was that this strategy also requires great effort, mostly because one has to extensively reword simple nouns which could be gendered with inclusive characters. For example, the word ``high school student'' is gender-specific in German and neutral alternatives are rather uncommon. Gender-inclusive characters were perceived as the easiest assignment by most of the participants (n=7) and relatively straightforward to use. Neosystems were never selected as the easiest strategy. Rather, they were the most difficult assignment (n=6) and considered particularly challenging. Nevertheless, the resulting cognitive overload did not always increase translation or PE times. No major differences emerged between the translation and PE tasks in the first and second assignments. When using neosystems, post-editors selected them more often as the most difficult strategy and also reported a more negative experience with them. Finally, half of the participants (TR=4, PE=2) commented that using only one strategy per text is limiting and they would rather use a mix of gender-neutral rewording and gender-inclusive characters (TR=3, PE=5). 

Table \ref{tab:interviews_GNR} provides an overview of the most interesting statements that participants made concerning gender-neutral rewording, while the interview transcripts in German are linked in the Appendix. While only P1 found it the most difficult strategy, half of the participants (TR=3, PE=3) commented that it requires great effort because, for example, it is sometimes difficult to find alternatives to gendered terms (TR=3, PE=3), which ``are otherwise totally natural for me, for example, `students' '' (P4). In general, participants commented that they had to restructure the texts and hence constantly reflect on possible alternatives. Nevertheless, this was the easiest assignment for five participants (TR=3, PE=2). There was some consensus regarding the fact that rewording is generally easy (TR=1, PE=4) because one can use one's own language repertoire without introducing new characters and/or forms. While some recognised that rewording requires creativity (TR=3, PE=1), this was considered both positive and negative. For example, P7 commented that they ``could think in every possible direction [...] I could remove personal references, and use proper names [...] I simply had great margin to reformulate''. Conversely, P2 felt that the margin for creativity was ``really unnatural'' as well as ``time-consuming and difficult''. 

\begin{table}[ht]
    \caption{Statements from the interviews regarding the first assignment}
    \centering
    %\addtolength{\leftskip}{-0.5cm}
    \begin{tabularx}{\textwidth}{QQrrr}
 \lsptoprule
        \textbf{Statement} & \textbf{Participant} & \textbf{TR} & \textbf{PE} & \textbf{Total} \\
\midrule
       It was the most difficult assignment & P1 & 1 & 0 & 1 \\
       It requires great effort & P1, P4, P6, P9, P10, P11 & 3 & 3 & 6 \\
       It was the easiest assignment & P3, P5, P6, P7, P12 & 3 & 2 & 5 \\
       It is generally easy to apply & P3, P8, P9, P10, P12 & 1 & 4 & 5 \\
       Creativity is allowed and/or required & P2, P5, P6, P7 & 3 & 1 & 4 \\
       It is difficult to find gender-neutral alternatives & P2, P4, P6, P9, P10, P11 & 3 & 3 & 6 \\
       I would prefer to use gender-inclusive characters & P4, P6, P9, P10 & 2 & 2 & 4 \\
       The translation reads differently than the source text & P1, P5, P6, P11 & 3 & 1 & 4 \\
       It is difficult in reference to a single person & P2, P5, P11 & 2 & 1 & 3 \\
       I am not satisfied with my solutions & P2, P6, P8, P11 & 2 & 2 & 4 \\
       Difficulties with ``actor and musician'' & P1, P2, P4, P6, P8, P10, P11 & 4 & 3 & 7 \\
       Difficulties with ``student'' & P4, P7, P8, P9, P10, P12 & 1 & 5 & 6 \\
       Difficulties with pronouns & P1, P2, P5, P10, P11 & 3 & 2 & 5 \\
       Difficulties with ``students'' & P6, P8, P9 & 1 & 2 & 3 \\
\lspbottomrule
    \end{tabularx}
    \label{tab:interviews_GNR}
\end{table}

Despite the generally positive attitudes towards this strategy, four participants (TR=2, PE=2) expressed a preference for gender-inclusive characters. In general, gender-neutral rewording is believed to shift the meaning of the text because the translation reads differently than the source text (TR=3, PE=1), for example, P11 stated ``I do not know whether there is some meaning loss [...] It is not something that Cal [the non-binary character in the text] would say, but rather something general''. This strategy is also considered difficult when referring to a single person (TR=2, PE=1) because it is normally used to neutralise plural references. As a consequence, some participants (TR=2, PE=2) stated that they were not satisfied with the solutions they found for the translation of gender-specific terms. The most common challenges in the first assignment were the translation of ``actor and musician'' (TR=4, PE=3), ``student'' (TR=1, PE=5), pronouns (TR=3, PE=2), and ``students'' (TR=1, PE=3). For instance, P8 stated that they had ``great difficulties with the terms `actor' and `musician'. I had to reflect for quite a long time on how I could reword them'' whereas P5 felt that ``(in the case of singular \textit{they}) I repeated the proper noun of the character too often [...] The general challenge is to avoid many repetitions so that the text is fluent but also [...] That the text is not too abstract'' (see, for example, P5's translation in the materials linked in the Appendix). These difficulties are also confirmed by the analysis of the screen activity tracking data, which was higher in the third segment containing the aforementioned nouns and in the last containing several pronouns.
%check for strategy limitation

As shown in Table \ref{tab:interviews_GIC}, even though two participants (TR=1, PR=2) commented that gender-inclusive characters were the most difficult assignment for them, there was quite a consensus that it is the easiest assignment (TR=3, PE=4) and generally easy to use (TR=1, PE=5). For example, P12 commented that ``[the non-binary person] was constantly mentioned in different inflections, dative, genitive (see the source texts and the MT outputs linked in the Appendix). And one could not actually avoid the pronouns'' while P10 thought that ``it is something we are already used to''. Most of the participants also agreed that adjective declension is straightforward (TR=3, PE=5), even though they have to combine masculine and feminine suffixes with a character if and when they differ. Only two participants (P1, P8) felt that the adjective declension is at times difficult. Also in this case, some participants (TR=1, PE=2) commented that characters are more difficult to use in reference to a single person, because this strategy is quite common in written language to avoid masculine generics. 

\begin{table}[ht]
    \caption{Statements from the interviews regarding the second assignment}
    \centering
    %\addtolength{\leftskip}{-0.5cm}
    \begin{tabularx}{\textwidth}{QQrrr}
    \lsptoprule
        \textbf{Statement} & \textbf{Participant} & \textbf{TR} & \textbf{PE} & \textbf{Total} \\
\midrule
       It was the most difficult assignment & P5, P12 & 1 & 1 & 2 \\
       It was the easiest assignment & P1, P2, P4, P8, P9, P10, P11 & 3 & 4 & 7 \\
       It is generally easy to apply & P1, P7, P8, P9, P10, P11 & 1 & 5 & 6 \\
       The adjective declension is straightforward & P4, P5, P6, P8, P9, P10, P11, P12 & 3 & 5 & 8 \\
       It is difficult to use in reference to a single person only & P5, P10, P12 & 1 & 2 & 3 \\
       It affects readability & P3, P4, P5, P6, P7, P8, P9, P11, P12 & 4 & 5 & 9 \\
       Effort to ensure consistent use & P1, P3, P4, P5, P6, P7, P11 & 5 & 2 & 7 \\
       Difficulties with ``doctor'' & P5, P6, P7, P8, P9, 11, P12 & 2 & 5 & 7 \\
       Uncertainties with ``recurring guest star'' & P4, P7, P8, P10, P11 & 1 & 4 & 5 \\
       Uncertainties how to combine forms & P2, P6, P8, P10, P11 & 2 & 3 & 5 \\
       Thought of masculine and feminine forms to use it & P1, P2, P3, P4, P5, P7, P8, P10, P12 & 5 & 4 & 9 \\
\lspbottomrule
    \end{tabularx}
    \label{tab:interviews_GIC}
\end{table}

In general, most of the translators (n=4) and post-editors (n=5) agreed that the frequent use of characters negatively impacts readability which could possibly transform the text into a ``storm of characters'' (P5). The greatest challenge concerning gender-specific terms was ``doctor'' (TR=2, PE=5) because, as already explained, the masculine and feminine form of the noun differ with an umlaut on the first letter, i.e. ``Arzt'' (male doctor)/``Ärztin'' (female doctor). This is uncommon; usually the word stem is not altered and the suffix \textit{-in} is added to the masculine form. Hence, participants were not sure how to use gender-inclusive characters. Another challenge that was often mentioned was the translation of terms like ``recurring guest star'' (TR=1, PE=5). The interviews made clear that the problem did not only concern finding an inclusive form. Sometimes it was difficult to find a viable translation in German or to understand whether gender-inclusive characters were to be used at all costs, as explained by P7: ``my question was `do I have to reformulate `member'? [...] This is neutral in German as well''.
Other issues include deciding how to combine masculine and feminine forms of articles and pronouns (e.g. ``der*die'' (masculine*feminine definite article) or ``die*der'' (feminine*masculine definite article)) (TR=2, PE=3) and ensure their consistent use as well as not to overlook gendered forms (TR=2, PE=5). For example, P1 looked at their translation and commented ``So here, too, the feminine comes first. Here too. Well, with the interrogative pronouns it is inevitably the same and with the nouns it is the other way round [...] It is not a very consistent strategy''. 

These uncertainties and difficulties are also confirmed by the observation protocols (see materials linked in the Appendix). First, the highest number of screen activities occurred in the first segment containing the noun ``doctor'', followed by the fourth segment with ``recurring guest star'', and ``recurring cast member''. Second, participants generally made more changes than in the first assignment and part of these included changing the order of masculine and feminine forms. Furthermore, this strategy does not seem to challenge the gender binary and may in fact reinforce it since the majority of the participants stated they needed to think of feminine and masculine forms of gender-specific terms in order to be able to combine them with characters (TR=5, PE=4). Finally, P5 made a very interesting remark: they translated the text first and subsequently adjusted gender references. 


Table \ref{tab:interviews_N} provides an overview of statements made concerning the use of neosystems. This strategy was generally considered the most difficult (TR=2, PE=4), and participants described it as challenging (TR=1, PE=3), e.g. ``it took up a lot of my brain capacity so that I often had the feeling there might be less capacity left for finding other nice formulations'' (P6). Nevertheless, for P2 and P11 the experience in the third assignment was better than they expected, e.g. ``yesterday I read the handout and thought that it sounded artificial [...] Today I found it actually pleasant'' (P2), and another three participants (P1, P10, P12) found it positive. 

\begin{table}[ht]
\caption{Statements from the interviews regarding the third assignment}
    \centering
        %\addtolength{\leftskip}{-0.5cm}
    \begin{tabularx}{\textwidth}{QQrrr}
    \lsptoprule
        \textbf{Statement} & \textbf{Participant} & \textbf{TR} & \textbf{PE} & \textbf{Total} \\
\midrule
      The experience was better than expected & P2, P11 & 1 & 1 & 2 \\
      It was a positive experience & P1, P10, P12 & 1 & 2 & 3 \\
      It was the most difficult assignment & P1, P3, P7, P9, P10, P11 & 2 & 4 & 6 \\
      It is challenging & P6, P8, P9 & 1 & 2 & 3 \\
      Training is required & P1, P3, P5, P6, P7, P8, P10, P11, P12 & 4 & 5 & 9 \\
      It feels like a foreign/artificial language & P3, P6, P8, P10, P11, P12 & 2 & 4 & 6 \\
      It affects readability & P8, P9, P10, P11, P12 & 0 & 5 & 5 \\
      Uncertainties about use & P3, P5, P9, P12 & 2 & 2 & 4 \\
      Uncertainties about correctness & P5, P6, P8, P9, P10, P11, P12 & 2 & 4 & 6 \\
      Difficulties with ``nanny'' & P2, P3, P4, P5, P6, P7, P8, P9, P11, P12 & 5 & 5 & 10 \\
      Difficulties with ``co-creators'' & P2, P5, P6, P7, P12 & 3 & 2 & 5 \\
      Difficulties with declensions & P1, P2, P3, P4, P5, P6, P8, P9 & 6 & 2 & 8 \\
      Need for support & P5, P8, P9, P10 & 1 & 3 & 4 \\
      Difficulties in selecting word class & P5, P10, P11 & 1 & 2 & 3 \\
      Resources used for the whole experiment & P1, P2, P7, P8, P9, P10, P11, P12 & 2 & 6 & 8 \\
      First content, then neosystem & P8, P11 & 0 & 2 & 2 \\
\lspbottomrule
    \end{tabularx}
    \label{tab:interviews_N}
\end{table}

There was a general consensus that further training would be required to effectively apply the strategy (TR=4, PE=5) and that neosystems are unknown, thus being perceived as a foreign and/or artificial language (TR=2, PE=6) and affecting the target text readability (PE=5). P3, for instance, commented that ``everything is completely new and it does not sound how one would speak or write''. Participants were thus unsure about how the neosystems are used even though they had access to a handout describing them (TR=2, PE=2). They were also unsure about whether they applied the selected system correctly (TR=2, PE=4). Specific challenges concerned the translation of ``nanny'' (TR=5, PE=5), because the profession is stereotypically feminine and the German equivalents are of feminine gender or have a different connotation. The term ``co-creators'' caused some difficulties as well. In this case, as in the previous assignment with ``recurring cast member'', this was sometimes due to terminology rather than GFL, e.g. ``it gave me a big headache. I then thought about what I would say in the non-gendered version'' (P6). Finally, participants struggled with declensions, especially adjectives (TR=6, PE=2) and stated that they used the handout provided or other resources during the whole task (TR=2, PE=6). They also commented that they would prefer to work with and/or ask another person to proofread their target text in order to ensure the correct use of the systems. Interestingly, P8 and P11 commented that their strategy was to first concentrate on the content and then post-edit all gender references with the selected neosystem.

While no clear-cut differences emerged between the translation and the PE group when commenting on the first and the second assignment, post-editors seemed to have more negative attitudes toward neosystems. For example, P7, P8, and P9 commented that they would refuse assignments with neosystems because they would require longer PE times and it would not be lucrative enough, e.g. ``I would have to put so much preparatory work into it that it probably would not be financially profitable'' (P8). Actually, throughout the interview with P8 (see transcript linked in the Appendix), a quite negative attitude towards more disruptive GFL strategies emerged. P9 also mentioned that they would otherwise increase their rates. Nevertheless, P7 needed less time to complete the third assignment than the first two, and P9 needed around 21 minutes for both the second and the third. While cognitively demanding, neosystems did not necessarily require more time and might therefore not be less lucrative. There were greater differences between the translation and the PE group when discussing the general use of MT as well as its (dis-)advantages, as shown in Table \ref{tab:interviews_MT}.

\begin{table}[ht]
    \caption{Statements from the interviews regarding MT}
    \centering
    %\addtolength{\leftskip}{-0.5cm}
    \begin{tabularx}{\textwidth}{QQrrr}
    \lsptoprule
        \textbf{Statement} & \textbf{Participant} & \textbf{TR} & \textbf{PE} & \textbf{Total} \\
\midrule
      I don't generally use MT & P1, P3, P6, P7, P10, P12 & 3 & 3 & 6 \\
      I use MT sometimes as a source of inspiration & P4, P7, P8, P9, P11 & 1 & 4 & 5 \\
      I would not use MT for similar texts & P1, P2, P3, P5, P9, P12 & 4 & 2 & 6 \\
      I would use MT for similar texts & P4, P11 & 1 & 1 & 2 \\
      MT is faster and/or cheaper & P6, P7, P8, P10, P11, P12 & 1 & 5 & 6 \\
      MT provides a translation draft & P4, P5, P7, P8, P10 & 2 & 3 & 5 \\
      It allows to concentrate on GFL & P10, P11 & 0 & 2 & 2 \\
      Extensive PE is required & P1, P2, P3, P5, P6, P7, P8, P9, P10, P12 & 5 & 5 & 10 \\
      MT lacks GFL & P1, P3, P4, P5, P6, P9 & 5 & 1 & 6 \\
      One is constrained by the MT draft & P4, P5, P7, P8, P10 & 2 & 3 & 5 \\
      It is easy to overlook gendered elements & P10 & 0 & 1 & 1 \\
\lspbottomrule
    \end{tabularx}
    \label{tab:interviews_MT}
\end{table}

Half of the participants (TR=3, PE=6) stated that they do not generally use MT. However, some (TR=1, PE=5), especially those in the PE group, use it as a source of inspiration, P11 said: ``So I like to work with machine translation, I use it mostly as a tool, I just have DeepL or context Reverso open most of the time''. In addition, most of the participants (TR=4, PE=2) find that MT is not a valuable tool for texts similar to those they translated or post-edited during the experiment, mostly because the MT outputs would require extensive PE of GFL since ``it is simply also a text or an area where you have to think a lot, and where so many people have to have a lot of background knowledge, a lot of thinking, simply reading between the lines'' (P1).
In comparison, two participants commented that MT can nevertheless be helpful.
The main advantages of MT concern the fact that PE is faster and/or cheaper than translation (TR=1, PE=5), especially because one has already a relatively good translation draft that can be adapted (TR=2, PE=3). Two participants even stated that the MT draft allows them to concentrate more on GFL since the output quality is rather good (PE=2). From the analysis of the (dis)advantages, it emerges that post-editors have a more positive attitude towards this technology than translators, who primarily remarked that MT outputs are biased and lack GFL (TR=5, PE=1) and extensive PE is required (TR=5, PE=5). Another disadvantage of this technology is that PE constrains creativity because the idea is to use as much as possible of the MT output and hence make fewer changes (TR=2, PE=3). Finally, one participant also commented that gendered elements can be overlooked. For instance, P10 stated that ``when I translate, I am in the text and I know which person this is about and whether they're binary or non-binary or whatever. And with machine translation, you really have to carefully check what elements refer to this person who is non-binary''.

Participants from the PE group were also asked about their specific impressions of the task. The majority (n=4) said it was difficult to decide to what extent the MT output should be adapted since the guidelines provided for full PE specified that one should use as much MT output as possible. Participants generally thought that, apart from lacking GFL, the biggest challenge was the stylistic component of the texts (n=3) -- for which the MT outputs were not appropriate, even though they were good and comprehensible (n=5). Two participants also said that the high-quality MT allowed them to concentrate on GFL, even though they felt that they then overlooked other important aspects of the PE process, e.g. ``what is in the source text? Does it correspond to the target text?'' (P11). 
Other comments included that GFL obviously requires extensive PE (n=1), thus there is a need to constantly revise the MT outputs to ensure that gender references are not overlooked (n=1). In general, participants (n=5) felt they achieved higher productivity through PE and they would have been slower if translating the study texts. Only one participant felt there would be no difference in productivity between the two tasks.

To summarise, strategies have different advantages and disadvantages. Gender-neutral rewording is considered a relatively easy approach to GFL but it requires creativity. While some participants felt free to experiment with the language, others felt it was unnatural, time-consuming, and challenging (see, e.g., P1's and P9's interview transcripts). Gender-inclusive characters were easier to use, mostly because people are already used to seeing them in newspapers or administrative texts. However, there is no standard concerning how to concatenate masculine and feminine forms with characters and some participants felt that this is arbitrary. In addition, the sole use of characters can negatively impact readability. Hence, participants expressed a preference for combining this strategy with gender-neutral rewording. Neosystems were the most challenging strategy mostly because they are unknown to the general public and require new grammatical rules, thus feeling like a foreign, unnatural language. Attitudes towards this strategy were more negative in the PE group: participants felt that this strategy negatively impacts their productivity, and so would at least require them to increase their rates. Participants still had a sceptical attitude towards MT; many of them do not use it, and find it unsuitable for creative and heavily gendered texts. However, post-editors found it increases their translation speeds and can be a valuable tool despite its biased outputs. 


%desires and real assignments


%----------------------------------------------------------------------------------------
%	STRATEGY SELECTION AND APPLICATION
%----------------------------------------------------------------------------------------

\section{Strategy selection and application} \label{sec:strategiesused}
%inserire nella metodologia quanto fare l'annotazione delle strategie sia un processo normativo che richiede di stabilire cosa sia sbagliato e/o non permesso siccome in realta' non ci sono degli standard ben definiti di come il linguaggio inclusivo si utilitzza
This section presents an overview of the strategies used by participants in each text, followed by an error analysis. In the first assignment with gender-neutral rewording, participants had to find creative ways to avoid gendered elements. This most often involved the use of gender-neutral words or rewording whole passages. In the second assignment with gender-inclusive characters, the majority of the translators and post-editors opted for gender star. However, this was put into practice quite differently. Furthermore, translators opted more often for gender-neutral words and rewording than post-editors, even if this was against the assignment instructions. Finally, the most selected neosystems were the Sylvain system and the Ens-forms. Interestingly, the former was a more popular choice in the PE group, whereas the translation group preferred the latter. No other relevant differences were found between the two groups. 

The error analysis also reveals that, apart from sporadic misgendering and use of masculine generics, participants rarely made mistakes in the first two assignments. Post-editors made fewer mistakes, but the difference from translators was minimal.  While using neosystems, there was a significant increase in errors due to difficulties in applying the systems. Misgendering was more common due to difficulties in finding an adequate German equivalent for the term ``nanny''. 

In the first text, nine gender phrases were selected for analysis and included references to non-binary individuals and mixed-gender groups. A total of 54 phrases were analysed and each was annotated for specific strategies used. The annotations totalled 55 and 59 for the translations and post-edited texts respectively, as shown in Table \ref{tab:strategies1}. This difference was due to more strategies being used for the same phrase. Gender-neutral words were the most common strategy in both groups (TR=22\% and PE=24\%), followed by rewording (TR=20\% and PE=22\%), and pronoun omission (TR=16\% and PE=12\%). Personal and possessive pronouns could sometimes be omitted and, when this was not possible, participants usually repeated the actor's or character's name (TR=7\% and PE=8\%). Another relatively common strategy (although only in the translations) was the use of participial forms. However, articles for these must be gendered in the singular form, e.g. in ``Schauspielende und Musizierende'' (actor and musician). For this reason, this strategy was and should be generally applied to plural nouns. While in the target text analysis other strategies were found less frequently, i.e. collective nouns, segment omissions and, against the instruction, gender-inclusive characters, sometimes no strategy was needed, i.e. in the case of pronouns referring to mixed-gender groups since the third-person plural pronoun in German is already gender-neutral. Table \ref{tab:examples1} shows examples of the great variety of strategies applied by participants in the first assignment.  

\begin{table}[!ht]
    \caption{Strategies used in the first assignment}
    \centering
    \begin{tabular}{lrrrr} \lsptoprule
        & \textbf{TR} & & \textbf{PE} &  \\ 
        \textbf{Strategy} & \textit{N} & \% & \textit{N} & \% \\
\midrule
        Gender-Neutral Words & 12 & 22 & 14 & 24 \\ %
        Rewording & 11 & 20 & 13 & 22 \\ %
        Pronoun Omission & 9 & 16 & 7 & 12 \\ %
        Participle & 7 & 13 & 1 & 2 \\ %
        Name Repetition & 4 & 7 & 5 & 8 \\ %
        Collective Noun & 1 & 2 & 2 & 3 \\ %
        Relative Pronoun & 1 & 2 & 1 & 2 \\ %
        Segment Omission & 0 & 0 & 2 & 3 \\ %
        Gender-Inclusive Character & 1 & 2 & 2 & 3 \\ %
        No Strategy & 9 & 16 & 12 & 20 \\ \midrule
        Total & 55 & 100 & 59 & 100 \\
\lspbottomrule
    \end{tabular}
    \label{tab:strategies1}
\end{table}


\begin{table}[!ht]
    \caption{Examples of strategies from the first assignment}
    \centering
    \begin{tabularx}{\textwidth}{QQQQ}
    \lsptoprule
        \textbf{Strategy} & \textbf{Source Text} & \textbf{TR} & \textbf{PE} \\
\midrule
        Gender-neutral words & actor and musician & Schauspielperson und Musikmensch & Schauspiel- und Musikprofi \\
        \tablevspace
        Rewording & it (the uniform) makes them feel more comfortable in who they are & weil sich diese angenehmer anfühlt & da dies dem Identitätsgefühl der nicht-binären Person am meisten entspricht \\
        \tablevspace
        Pronoun omission & their right & das Recht & das Recht \\
        \tablevspace
        Name repetition & They refuse to wear & Cal weigert sich & Cal weigert sich \\
        \tablevspace
        No strategy & shining a bright light on those who identify as non-binary & setzt den Fokus auf jene, die sich als nicht-binär identifizieren & stellt diejenigen, die sich als nicht-binär identifizieren, in den Mittelpunkt \\
\lspbottomrule
    \end{tabularx}
    \label{tab:examples1}
\end{table}

In the second text, 12 gender phrases were selected for analysis, and 81 and 77 annotations were made to the translations and post-edited texts respectively. In this case as well, more strategies were used for a specific phrase, showing that participants sometimes mixed gender-inclusive characters with rewording. Reasons for this include that some neutral words, e.g. ``Publikum'' for ``audiences'', are relatively common and hence a natural translation choice while rewording should ensure better readability because it helps reduce the number of characters used to concatenate masculine and feminine forms. As in the previous assignment, sometimes no specific GFL strategy was needed, e.g. when ``they'' was used as a third-person singular pronoun in the source text. A general overview of the applied strategies is presented in Table \ref{tab:strategies2}.

\begin{table}[!ht]
    \centering
    \caption{Strategies used in the second assignment}
    \begin{tabular}{lrrrr} \lsptoprule
        & \textbf{TR} & & \textbf{PE} &  \\ 
        \textbf{Strategy} & \textit{N} & \% & \textit{N} & \% \\
\midrule
         Gender-Inclusive Characters & 39 & 48 & 44 & 59 \\ %
         Gender-Neutral Rewording & 28 & 35 & 17 & 23 \\ %
         No Strategy & 14 & 17 & 13 & 18 \\ \midrule
         Total & 81 & 100 & 77 & 100 \\
\lspbottomrule
    \end{tabular}
    \label{tab:strategies2}
\end{table}

Interestingly, participants in the PE group used gender-inclusive characters more often (59\%) than those in the translation group (48\%) who used gender-neutral words and/or rewording more frequently (TR: 35\%, PE: 23\%). More interesting, however, is to analyse which gender-inclusive characters were selected by the participants and how they were applied. 

Gender star (*) was the most selected character, namely by eight participants (TR=3, PE=5), followed by underscore (\_) (TR=2), and colon (:) (TR=1, PE=1). Variations were found in their use. Here the focus is on gender star because it was applied in such a diverse way. While nouns are created with the word stem or the masculine form adding the gender-inclusive character and then the feminine suffix ``-in'', there is more freedom when using articles and pronouns. The following variations were found:
\begin{itemize}
\setlength\itemsep{-0.5em}
    \item masculine forms were always used first, e.g. ``der*die Schauspieler*in'' (masculine*feminine article + masculine*feminine noun) (n=2).
    \item feminine forms were used first in articles and pronouns, e.g. ``die*der Schauspieler*in'' (feminine*masculine article + masculine*feminine noun) (n=4).
    \item Pronouns and articles were combined in a new form, e.g. ``die*r Schauspieler*in'' (n=1). This is, however, not possible for possessive pronouns.
    \item Slash (/) was used to combine masculine and feminine forms of pronouns and articles, e.g. ``der/die Schauspieler*in'' (n=1). This variation is rather uncommon. 
\end{itemize}

In the third text, ten gender phrases were selected for analysis. A total of 63 and 65 annotations were made to the translations and post-edited texts respectively. In this case as well, some phrases were annotated twice because participants used more than one strategy. An overview is provided in Table \ref{tab:strategies3}. In comparison to the second assignment, no differences emerge between translation and PE: all participants mostly used neosystems and less frequently resorted to gender-neutral rewording. 

\begin{table}[!ht]
    \centering
    \caption{Strategies used in the third assignment}
    \begin{tabular}{lrrrr} \lsptoprule
        & \textbf{TR} & & \textbf{PE} &  \\ 
        \textbf{Strategy} & \textit{N} & \% & \textit{N} & \% \\ \midrule
         Neosystems & 45 & 71 & 46 & 70 \\ %
         Gender-Neutral Rewording & 14 & 22 & 13 & 20 \\ %
         No Strategy & 4 & 6 & 6 & 9 \\ \midrule
         Total & 63 & 100 & 65 & 100 \\
\lspbottomrule
    \end{tabular}
    \label{tab:strategies3}
\end{table}

However, differences were found in the system selection:
\begin{itemize}
\setlength\itemsep{-0.5em}
    \item The Ens-forms, e.g. ``Ens spielt dens Nanny Sabi Mehboob'' (they star as nanny Sabi Mehboob), were selected by 5 participants (TR=4, PE=1).
    \item The Sylvain system, e.g. ``Dabei handelt sich um eine Komödie über einir genderfluidin Nanny'' (this is a comedy about a gender-fluid nanny) was selected by another 5 participants (TR=1, PE=4).
    \item The NoNa-System, e.g. ``Baig ist eint südasiatische, queere, muslimische Schauspieler*in'' ( Baig is a South Asian, queer, Muslim actor) was selected by one participant (PE=1).
    \item One participant selected a system which they use in their everyday life and consists of neopronoun \textit{dey/deren}, the gender-fair suffix \textit{-i}, and the use of neuter for articles and adjective declension, e.g. ``Eine Komödie über ein genderfluides Kindermenschi'' (a comedy about a gender-fluid nanny).
\end{itemize}
While three participants stated that they selected the system arbitrarily, four used the Ens-forms because they considered it the easiest strategy and three chose the Sylvain system as it is the most developed, with the introduction of a fourth gender in the German grammar. Interestingly, translators opted for easier neosystems, while post-editors for more sophisticated ones. As explained in Chapter \ref{Discussion}, PE perhaps requires less cognitive load, allowing for greater focus on neosystems.


Errors in each target text segment were also annotated. Errors included misgendering, misrecognition of singular \textit{they}, i.e. translation in the third-person plural German pronoun, masculine generics, use of false terms, and mistakes in the application of the strategy, e.g. use of wrong pronouns or articles. Masculine generics are not an error \textit{per se}, but they represent a sexist use of language. An overview of the results is presented in Table \ref{tab:error}, where the number of errors for each analysed segment is reported. As already explained, the number of analysed segments refers to the gender phrases that required GFL in translation and in PE. For each assignment (1, 2, 3), results for translation and PE are both reported in terms of absolute (N) and relative (\%) frequencies. In Table \ref{tab:error2}, the number of participants who made the errors is reported. This aims to provide an overview of how common certain error types were among the participants.

\begin{table}[!ht]
    \centering
    \caption{Overview of the error analysis}
    \resizebox{\textwidth}{!}{\begin{tabular}{lrrrrrrrrrrrr} \lsptoprule
      \textbf{Assignment}       & 1           &    &             &    & 2           &    &             &    & 3           &    &             & \\
                                & \textbf{TR} &    & \textbf{PE} &    & \textbf{TR} &    & \textbf{PE} &    & \textbf{TR} &    & \textbf{PE} & \\
      \textit{Error Type}       & \textit{N/54}  & \% & \textit{N/54}  & \% & \textit{N/72}  & \% & \textit{N/72}  & \% & \textit{N/60}  & \% & \textit{N/60}  & \% \\
\midrule
      Misgendering              & 2           & 4  & 1           &  2 & 1           & 1  & 0            & 0   & 7           & 12  & 7           & 12 \\
      Misrecognition Sing. \textit{They} & 1             & 2   &    0         &  0  & 4           & 6  & 0            & 0   &    0         & 0    &  0           & 0 \\
      Masculine Generics             &  0           & 0   & 1           &  2 & 0            & 0   & 2           & 3  & 0            & 0   &    0         & 0 \\
      False Term                & 6           & 11 &  0           &  0  & 0            & 0   &  0           &  0  & 0            & 0   & 0            & 0 \\
      Mistake                   & 0            & 0   & 0            & 0    & 5           & 7  & 1           & 1  & 17          & 28 & 17          & 28 \\
      Total                     & 9        & 17 & 2        &  4 & 10       & 14 & 3       & 4  & 24       & 40 & 24          & 40 \\
\lspbottomrule
    \end{tabular}}
    \label{tab:error}
\end{table}

In general, the number of errors was quite low in the first (TR=17\%, PE=4\%) and in the second assignment (TR=14\%, PE=3\%); it increased in the third (TR=40\%, PE=40\%). The most common error in the first segment was the use of ``Studierende'' (university students) as the German equivalent for ``students'' which actually referred to high-school in the text. This mistake was not found in the post-edited texts, most probably because the MT selected the appropriate translation. Misgendering was found twice in the translations and once in the post-edited text. P3 used the participial form ``Studierende'' in the singular to refer to the non-binary character mentioned in the text. Participial forms are neutral only in the plural, otherwise a masculine or feminine article must be used. A similar mistake was made by P10 when post-editing ``actor and musician''. P5 misgendered the actor mentioned in the text using masculine forms. During the interview, they commented that they thought the adjective ``non-binary'' made the actor's gender clear and did not know how to translate ``actor and musician'' by rewording. Finally, P3 also misrecognised singular \textit{they} in one passage: the phrase ``They refuse to wear'' became ``Die Studierenden weigern sich [...] zu tragen'' (The students refuse to wear). 

The most common mistake in the second text was the misrecognition of singular \textit{they}. This, however, was due to the ambiguity of the text. In the final passage, ABC's character description of the non-binary doctor portrayed in the TV series \textit{Grey's Anatomy} is reported but another character is also mentioned. Therefore, it was difficult to understand whether ``they'' referred to the non-binary doctor or both characters. This aspect was also mentioned by participants in the interviews: while the use of the said pronoun is generally clear (mentioned by five participants), certain text passages were ambiguous (mentioned by eight participants). Another common error type concerned the use of gender-inclusive characters, mostly in the case of ``doctor'' for which three participants (TR=2, PE=1) concatenated the masculine and the feminine form of the noun (e.g. ``Arzt:Ärztin'') because these differ not only for the suffix but also for the umlaut on the first letter. Another participant in the translation group made a mistake in the adjective declension. Finally, misgendering was found in one text only: P1 used a masculine article in front of the title ``Dr.'' referring to the non-binary character mentioned in the text.

\begin{table}[!ht]
    \centering
    \caption{Number of participants who made errors}
    \begin{tabular}{lcccccc} \lsptoprule
      \textbf{Assignment}       & 1              &                & 2              &                & 3           &              \\
\midrule
       \textit{Error Type}      & \textbf{TR}    & \textbf{PE}    & \textbf{TR}    & \textbf{PE}    & \textbf{TR} & \textbf{PE} \\
      Misgendering              & 2              & 1              & 1              & 0              & 3           & 3           \\
      Misrecognition Sing. \textit{They} & 1              & 0              & 2              & 0              & 0           &  0          \\
      Masculine Generics             & 0              & 1              & 0              & 1              & 0           &    0       \\
      False Term                & 3              & 0              & 0              & 0              & 0           & 0          \\
      Mistake                   & 0              & 0              & 3              & 1              & 5          & 6          \\
\lspbottomrule

    \end{tabular}
    \label{tab:error2}
\end{table}

In the third text, misgendering was quite common because participants faced difficulties in the translation of ``nanny''. While some opted for gender-fair versions of the term ``Kinderbetreuer'' (caregiver), others did not because they found its connotations are different. ``Nanny'' has a more colloquial and familiar tone so that one can maybe picture a young person trying to earn some money while also studying. ``Caregiver'' would be more formal and might convey the idea of a professional specialised in childcare. As a consequence, half of the participants (TR=3, PE=3) kept the term in the target texts and applied GFL to articles and/or pronouns, as in ``Nin spielt din Nanny'' (they star as nanny). The number of errors in the use of the neosystems was very high -- 28\% of the analysed segments both in translation and PE contained mistakes. This also means that every participant in the study made at least one mistake except for P4. In fact, they inadvertently used false pronouns throughout the text (e.g. ``sein'', his/its) but these errors were counted as misgendering only. During the interview, the participant realised the false use of pronouns and thus commented on their translation adding the correct alternatives.

The increased number of mistakes in the third assignment confirms the intensified cognitive effort commented on by participants during the post-study interviews as well as the higher intensity of screen activity observed during the translation and PE process. Participants were not sure about the (correct) use of neosystems and constantly needed resources in order to apply them. Hence, they modified their target texts and looked for resources more often than in the first two assignments. Minimal differences were found between translation and PE in terms of the total number of errors (see Table \ref{tab:error}). Since a large number of mistakes made by one participant could affect the values for absolute and relative occurrences among the analysed gender phrase, Table \ref{tab:error2} reports only the number of participants who made mistakes. In the first two assignments, no post-editor misrecognised singular \textit{they}, while one translator in the first and two translators in the second assignment did so. Besides, participants in the PE group never used the false term ``Studierende'' (university student), whereas half the translation group did. 

To conclude, participants showed great levels of creativity in the first assignment since they had to avoid gender references at all costs. They usually opted for gender-neutral words and rewording whole passages, often also changing perspectives or abstracting information, e.g. a translation for ``it (the uniform) makes them (the non-binary person) feel more comfortable'' was ``weil sich diese angenehmer  fühlt'' (because this feels better). In the second assignment, participants largely opted for gender star but its implementation varied. As a matter of fact, there is no established norm for the use of gender-inclusive characters and some participants (n=4) expressed the wish to have GFL dictionaries and/or more literature on the topic to make better-informed decisions. In the third assignment, participants from the translation group opted for easier systems while those from the PE group for more sophisticated strategies. While in the first two texts the number of mistakes was limited with post-editors performing slightly better than translators, it increased in the third text. The highest number of mistakes concerned the application of specific GFL strategies, whereas misgendering was found quite rarely except for the third text in which the translation of the term ``nanny'' proved particularly challenging.


%----------------------------------------------------------------------------------------
%	BIAS IN THE MT OUTPUTS
%----------------------------------------------------------------------------------------
\section{Gender bias in the MT outputs}
To contextualise the analysis of gender bias in DeepL's MT output, this section starts with a brief overview of the system's underlying architecture and functionality. DeepL Translator is an NMT system introduced in 2017 by the German company DeepL GmbH. It is based on proprietary developments derived from the Transformer architecture introduced by \textcite{vaswani2017attention}. Although DeepL's full model architecture has not been publicly released, available information and independent analyses suggest that it adheres to the encoder-decoder structure characteristic of Transformer-based models. These models rely on self-attention mechanisms, which enable dynamic contextual weighting of input tokens and are crucial for modelling long-distance dependencies in language \citep{bahdanau2015neural,vaswani2017attention,popel2020transforming}. This architecture is particularly advantageous for languages with rich morphology, flexible word order, or complex agreement systems -- factors closely linked to the manifestation of grammatical gender.

Like most NMT systems, DeepL’s default mode of operation is sentence-based. Translation units are segmented primarily at punctuation or line breaks, with each sentence processed independently \citep{laubli-etal-2018-machine,deepl_api_reference}. This sentence-level granularity offers computational efficiency but imposes significant limitations when modeling features that depend on discourse-level information, such as coreference resolution and gender agreement across sentences \citep{hardmeier2014discourse,voita2018context}. In the context of gender translation, this segmentation often leads to inconsistent gender assignments across contiguous discourse, especially when the referent’s gender is not lexically marked in the source language (e.g., English \textit{doctor} vs. French \textit{docteur/docteure}).

DeepL attempts to mitigate some of these limitations by offering contextual enhancement parameters through its API \citep{deepl_api_reference}. Specifically, the \texttt{split\_sentences} parameter allows users to disable automatic sentence segmentation, thereby enabling longer text segments to be translated as a whole. However, this setting risks incorrect chunking if punctuation is non-standard or ambiguous, and excessively long input strings can trigger truncation or decreased translation quality.

Additionally, DeepL provides a \texttt{context} parameter that allows for the inclusion of auxiliary text intended to inform the translation of the primary input \citep{deepl_context_guidance}. According to DeepL’s API documentation, this feature is particularly beneficial for ``short, low-context source texts'' and it enables the system to use adjacent information without directly translating it. Nonetheless, this parameter operates only within the scope of a single API call and does not permit discourse continuity across multiple \texttt{text} parameters. As such, while the \texttt{context} parameter can improve local coherence, it does not solve the broader issue of document-level gender consistency unless the entire document is submitted as a single text unit.

In sum, while DeepL's Transformer-based architecture offers notable advantages in handling syntactic complexity and producing fluent output, its sentence-level processing paradigm imposes inherent limitations for discourse-level phenomena such as gender coherence. These constraints are particularly pertinent for research into gender bias in NMT, as gender assignments often require information from broader discourse context, which may be inaccessible to sentence-level translation models.


In Section \ref{sec:technicaleffort}, it was highlighted that the main difference between the translation and PE process lies in the number of changes: post-editors were presented with MT outputs in which the gender references were generally wrong. MT phrases containing gender references to be post-edited were annotated for three typologies of gender bias, namely
(i) misgendering,
(ii) mistranslations of singular \textit{they},
(iii) masculine generics. Table \ref{tab:MTerrors} provides an overview of the results of the gender bias analysis for each machine-translated text. Note that one phrase could contain different typologies of bias, thus the number of gender-biased phrases and the total number of annotations in Table \ref{tab:MTerrors} do not always match.

\begin{table}[ht]
\centering
\small
%\addtolength{\leftskip}{-1.5cm}
    \caption{Bias in the MT outputs}
    \begin{tabularx}{\textwidth}{lp{2cm}Q}
    \lsptoprule
        \textbf{Text No.} & \textbf{MT errors per phrase} & \textbf{Types of MT errors} \\
\midrule
        1 & 6/9 & Misgendering (2), mistranslations of singular \textit{they} (3), masculine generics (1)\\ %
        2 & 10/12 & Misgendering (5), mistranslations of singular \textit{they} (4), masculine generics (2)\\ %
        3 & 9/10 & Misgendering (6), mistranslations of singular \textit{they} (3), masculine generics (1) \\
\lspbottomrule
    \end{tabularx}
    \label{tab:MTerrors}
\end{table}

Six out of the nine analysed gender references in the first text presented gender bias in the MT outputs. The other three phrases contained a plural indefinite pronoun, i.e. ``those'', and plural \textit{they} in the possessive case, i.e. ``their right'' and ``their identity''. When plural nouns were translated, masculine generics were used, e.g. ``students'' as ``Schüler''. Nouns in reference to the non-binary actor and character mentioned in the text were translated with masculine forms whereas singular \textit{they} was not recognised as such and translated with the third-person plural pronoun ``sie''.

Similarly, ten out of the 12 analysed gender references in the second text were biased as well. The character depicted in the text was always misgendered except for two phrases. For the first, i.e. ``Fightmaster was promoted from recurring guest star to a recurring cast member'', the German equivalent is gender-neutral as well, i.e. ``Fightmaster wurde von einem wiederkehrenden Gaststar zu einem wiederkehrenden Cast-Mitglied befördert''. The second contained the English title ``Dr.'' which is also sometimes used in German as gender-neutral. Otherwise, the actor and the character mentioned in the text were systematically misgendered and singular \textit{they} translated in the plural form. Interestingly, this time the term ``neuroscientist'' was translated in the feminine form, i.e. ``Neurowissenschaftlerin''. 

The same error typologies were found in the third text as well and only one phrase was translated correctly, i.e. ``audiences'' into the gender-neutral ``Publikum''.  The non-binary person mentioned in the text was generally misgendered with masculine, feminine and neuter forms. ``Nanny'' was translated as ``Kindermädchen'' which is grammatically neuter in German but socially female.  During the annotation, two MT errors were found that deserve further attention. First, there was a coreference mistake. In the English language text, “audiences” are mentioned in the first sentence, and “they” is used in the second as an anaphoric pronoun. In the German MT, the plural noun was translated into the singular, i.e. “das Publikum” (the public), and the third-person singular pronoun “sie” was used as an anaphor instead of the grammatically correct “es” (it). This might happen because of DeepL operating at the sentence level. Second, the series is about a nanny who is genderfluid. This identity term was translated into “geschlechtsspezifisch” (gender-specific) which is wrong.

In line with \textcite{prates2019assessing,piazzollagood,lardelli-etal-2024-building}, the analysis of the MT outputs seems to indicate that automated translation systems are still largely biased. A tendency to overuse masculine forms is noted while singular \textit{they} is translated in the plural form. Nevertheless, as explained in Section \ref{sec:cognitiveeffort}, post-editors found that the MT output increased their productivity, which is also confirmed by comparing translation and PE times. Sometimes, MT even allowed more concentration on the gender aspect of the task.
